% Originally: the \section{Solutions} of content/week-02.tex and content/week-03.tex

\section{Solutions}
\label{sec:completeness_solutions}

\begin{exercisesolution}[Proof of \cref{ex:cauchy_elementary_facts}]
% Generator: Claude Sonnet 5 (High)
\textbf{(a) Bounded.} Take $\varepsilon=1$: there is $N$ with $d(x_n,x_m)<1$ for all
$n,m\geq N$; in particular $d(x_n,x_N)<1$ for $n\geq N$. Setting
$M := \max\{d(x_0,x_N),\dots,d(x_{N-1},x_N),\,1\}$, the triangle inequality gives
$d(x_n,x_0) \leq d(x_n,x_N)+d(x_N,x_0) \leq 1+M$ for every $n$, so $(x_n)_n \subseteq B_{1+M}(x_0)$.

\textbf{(b) Convergent $\Rightarrow$ Cauchy.} If $x_n \to x$, fix $\varepsilon>0$ and choose $N$
with $d(x_n,x)<\varepsilon/2$ for $n\geq N$. Then for $n,m\geq N$,
$d(x_n,x_m) \leq d(x_n,x)+d(x,x_m) < \varepsilon$.

\textbf{(c) Cauchy + convergent subsequence $\Rightarrow$ convergent.} Let $(x_n)$ be Cauchy and
$x_{n_k} \to x$. The \qt{only if} direction is part \textbf{(b)} applied to the whole sequence.
For \qt{if}: fix $\varepsilon>0$. Cauchyness gives $N$ with $d(x_n,x_m)<\varepsilon/2$ for
$n,m\geq N$; convergence of the subsequence gives $K$ with $d(x_{n_k},x)<\varepsilon/2$ for
$k\geq K$. Since $n_k \to \infty$, pick $k\geq K$ with $n_k \geq N$. Then for every $n\geq N$,
\[d(x_n,x) \leq d(x_n,x_{n_k}) + d(x_{n_k},x) < \tfrac{\varepsilon}{2}+\tfrac{\varepsilon}{2} = \varepsilon,\]
so $x_n \to x$.
\end{exercisesolution}

\begin{exercisesolution}[Solution to \cref{ex:completeness_not_topological}]
% Generator: Claude Opus 5 (Medium)
Write $\varphi : (0,1] \to [1,\infty)$, $\varphi(x) := 1/x$. It is a bijection, and by
construction $d(x,y) = |\varphi(x)-\varphi(y)|$. Consequently, $d$ is nothing but the Euclidean
metric of $[1,\infty)$ transported back to $(0,1]$ along $\varphi$. Every part follows from that.

\begin{enumerate}[label=\textbf{(\alph*)}]
  \item Symmetry and the triangle inequality are inherited from $|\cdot|$ directly. For
    definiteness, $d(x,y)=0$ forces $\varphi(x)=\varphi(y)$, and $\varphi$ is injective, so
    $x=y$.
  \item The sequence $x_n := \tfrac1n$ lies in $(0,1]$ and is Cauchy for
    $d_{\text{eucl}}$, being Cauchy in $\mathbb{R}$. Its limit in $\mathbb{R}$ is $0 \notin X$,
    and limits in $\mathbb{R}$ are unique (\cref{prop:uniqueness_of_limits}), so it has no limit
    in $X$.
  \item Let $(x_n)$ be $d$-Cauchy. Then $\big(\varphi(x_n)\big)$ is Cauchy in $\mathbb{R}$,
    hence converges to some $\ell \in \mathbb{R}$. Since $\varphi(x_n) \geq 1$ for every $n$,
    also $\ell \geq 1$; in particular $\ell \neq 0$, so $x := 1/\ell$ lies in $(0,1]$. Then
    \[d(x_n,x) = \big|\varphi(x_n) - \varphi(x)\big| = \big|\varphi(x_n)-\ell\big|
      \longrightarrow 0,\]
    so $x_n \to x$ in $(X,d)$. Note where the boundedness of $X$ entered: it is what forces
    $\ell \geq 1$ rather than merely $\ell \geq 0$.
  \item $\varphi$ and $\varphi^{-1}(t) = 1/t$ are continuous, so $\varphi$ is a
    homeomorphism. A set is $d$-open exactly when its $\varphi$-image is open in $[1,\infty)$,
    exactly when it is $d_{\text{eucl}}$-open. Concretely, the same two-sided-ball argument as in
    \Cref{ex:heine_borel_fails}\textbf{(b)} works: on $(0,1]$ one has
    $|x-y| \leq d(x,y) \leq \tfrac{1}{x y}\,|x-y|$ with $\tfrac{1}{xy}$ locally bounded.
\end{enumerate}

\begin{remark}
So $d_{\text{eucl}}$ and $d$ have the \emph{same open sets}, the same convergent sequences, the
same continuous functions and the same compact sets, and yet one is complete and the other is
not. \textbf{Completeness is a property of the metric, not of the topology it generates.}

This completes a trio worth holding together, all three proved by the same device of comparing
two metrics that induce one topology:
\begin{itemize}
  \item \textbf{Compactness is topological.} It is defined by open covers alone
    (\cref{def:compact_metric_space}), so it cannot distinguish two metrics with the same open
    sets.
  \item \textbf{Boundedness is not:} see \cref{rem:boundedness_not_topological}.
  \item \textbf{Completeness is not:} this exercise.
\end{itemize}
That is exactly why \qt{compact if and only if closed and bounded} needs the standard metric, and
why the correct general replacement is \qt{complete and totally bounded}
(\cref{rem:compact_iff_complete_totally_bounded}). Both ingredients of the replacement are
non-topological, and only their combination is.
\end{remark}
\end{exercisesolution}

\begin{exercisesolution}[Proof of \cref{ex:completion_well_defined}]
% Generator: Claude Sonnet 5 (High)
\textbf{$\sim$ is an equivalence relation on $\mathcal{C}_X$.} Reflexivity is
$d(x_n,x_n) = 0 \to 0$. Symmetry is symmetry of $d$. For transitivity, if
$(x_n)\sim(y_n)$ and $(y_n)\sim(z_n)$, the triangle inequality gives
$0 \leq d(x_n,z_n) \leq d(x_n,y_n)+d(y_n,z_n) \to 0+0 = 0$, so $(x_n)\sim(z_n)$.

\textbf{$\overline d$ is well-defined.} First, the limit defining $\overline d$ exists: for
$(x_n),(y_n) \in \mathcal{C}_X$, the reverse triangle inequality gives
\[\big|d(x_n,y_n) - d(x_m,y_m)\big| \leq d(x_n,x_m) + d(y_n,y_m),\]
and the right-hand side $\to 0$ as $n,m\to\infty$ since $(x_n)$ and $(y_n)$ are each Cauchy, so
$\big(d(x_n,y_n)\big)_n$ is a Cauchy sequence of \emph{real numbers}, hence convergent. Note that
this is the completeness axiom for $\mathbb{R}$ itself from \emph{Analysis I}, and not an instance
of \cref{def:complete_metric_space} being proved circularly. Second, the limit
does not depend on the choice of representative: if $(x_n)\sim(x_n')$ and $(y_n)\sim(y_n')$, the
same reverse triangle inequality gives
$\big|d(x_n,y_n)-d(x_n',y_n')\big| \leq d(x_n,x_n')+d(y_n,y_n') \to 0$, so the two limits agree.

\textbf{$\overline d$ is a metric.} Fix representatives $(x_n)\in\overline x$, $(y_n)\in\overline
y$, $(z_n)\in\overline z$.
\begin{itemize}
  \item \textbf{Definiteness.} $\overline d(\overline x,\overline y) = 0$ means
    $\lim_n d(x_n,y_n) = 0$, i.e.\ exactly $(x_n)\sim(y_n)$, i.e.\ $\overline x = \overline y$.
  \item \textbf{Symmetry.} Immediate from $d(x_n,y_n)=d(y_n,x_n)$ for every $n$.
  \item \textbf{Triangle inequality.} Take $n\to\infty$ in
    $d(x_n,z_n) \leq d(x_n,y_n) + d(y_n,z_n)$, using that limits preserve weak inequalities.
\end{itemize}
Finally $d(x,y) = \overline d(\iota(x),\iota(y))$ holds because the constant sequences
$(x,x,\dots)$ and $(y,y,\dots)$ give $d(x_n,y_n) = d(x,y)$ for every $n$, so the limit is
$d(x,y)$ itself; and this identity makes $\iota$ injective, since $d(x,y)=0 \iff x=y$.
\end{exercisesolution}

\begin{exercisesolution}[Proof of \cref{ex:completion_is_complete}]
% Generator: Claude Sonnet 5 (High)
Let $(\xi_m)_{m=0}^\infty$ be Cauchy in $\overline X$, with $\xi_m = [(x_{m,n})_{n=0}^\infty]$ for
some Cauchy sequence $(x_{m,n})_n$ in $X$ representing $\xi_m$. Replacing $(x_{m,n})_n$ by a tail
of itself if necessary, which is harmless because a Cauchy sequence and any of its tails represent
the same class, differing in only finitely many terms, we may assume each representative is already
\newterm{regularized}:
\[d(x_{m,n},x_{m,n'}) < \tfrac1m \qquad \text{for all } n,n' \geq 0.\]
Set $y_m := x_{m,m}$, the diagonal sequence the exercise asks about. Letting $n'\to\infty$ in the
regularization bound (with $n=m$) gives
\begin{equation}
\overline d\big(\iota(y_m),\,\xi_m\big) = \lim_{n\to\infty} d(x_{m,m},x_{m,n}) \ \leq\ \tfrac1m.
\label{eq:completion_diagonal_close}
\end{equation}

\textbf{$(y_m)_m$ is Cauchy in $X$.} By the triangle inequality for $\overline d$ (already
established in \cref{ex:completion_well_defined}) and \eqref{eq:completion_diagonal_close},
\[d(y_m,y_{m'}) = \overline d\big(\iota(y_m),\iota(y_{m'})\big)
  \ \leq\ \overline d(\iota(y_m),\xi_m) + \overline d(\xi_m,\xi_{m'}) + \overline d(\xi_{m'},\iota(y_{m'}))
  \ \leq\ \tfrac1m + \overline d(\xi_m,\xi_{m'}) + \tfrac1{m'}.\]
Given $\varepsilon>0$, Cauchyness of $(\xi_m)_m$ gives $M_1$ with
$\overline d(\xi_m,\xi_{m'}) < \varepsilon/3$ for $m,m'\geq M_1$; taking also
$M \geq \max\{M_1,\,3/\varepsilon\}$ makes all three terms $< \varepsilon/3$ for $m,m'\geq M$, so
$(y_m)_m$ is Cauchy in $X$ and defines a class $\overline x := [(y_m)_m] \in \overline X$.

\textbf{$\xi_m \to \overline x$.} Using \eqref{eq:completion_diagonal_close} again and the same
triangle inequality,
\[\overline d(\xi_m,\overline x) \ \leq\ \overline d(\xi_m,\iota(y_m)) + \overline d(\iota(y_m),\overline x)
  \ \leq\ \tfrac1m + \lim_{k\to\infty}d(y_m,y_k).\]
The first term $\to 0$ as $m\to\infty$ by construction, and the second does too, being (for each
fixed $m$) bounded by $\sup_{k\geq m}d(y_m,y_k)$, which $\to0$ as $m\to\infty$ since $(y_m)_m$ is
Cauchy. Hence $\overline d(\xi_m,\overline x) \to 0$, i.e.\ $\xi_m \to \overline x$ in
$\overline X$, as required.
\end{exercisesolution}

\begin{exercisesolution}[Proof of \cref{ex:completion_universal_property}]
% Generator: Claude Sonnet 5 (High)
\textbf{Construction.} For $\overline x = [(x_n)_n] \in \overline X$, the sequence $(f(x_n))_n$
is Cauchy in $Y$, since $d_Y(f(x_n),f(x_m)) = d(x_n,x_m) \to 0$ as $(x_n)_n$ is Cauchy in $X$.
As $Y$ is complete, $f(x_n) \to y$ for a unique $y \in Y$
(\cref{prop:uniqueness_of_limits}); define $\overline f(\overline x) := y$.

\textbf{Well-defined.} If $(x_n)\sim(x_n')$, then $d_Y(f(x_n),f(x_n')) = d(x_n,x_n') \to 0$, so
$(f(x_n))_n$ and $(f(x_n'))_n$ have the same limit in $Y$ by uniqueness of limits; hence
$\overline f(\overline x)$ does not depend on the representative chosen.

\textbf{$f = \overline f \circ \iota_X$.} For $x\in X$, $\iota_X(x)$ is represented by the
constant sequence $(x,x,\dots)$, so $\overline f(\iota_X(x)) = \lim_n f(x) = f(x)$.

\textbf{$\overline f$ satisfies the displayed identity.} For $\overline x = [(x_n)_n]$,
$\overline y = [(y_n)_n]$, continuity of $d_Y$ in both arguments (itself $1$-Lipschitz, by the
same reverse-triangle-inequality argument as in \cref{ex:completion_well_defined}) gives
\[d_Y\big(\overline f(\overline x),\overline f(\overline y)\big)
  = d_Y\Big(\lim_n f(x_n),\ \lim_n f(y_n)\Big)
  = \lim_n d_Y\big(f(x_n),f(y_n)\big)
  = \lim_n d(x_n,y_n) = \overline d(\overline x,\overline y).\]

\textbf{Uniqueness.} Let $g : \overline X \to Y$ satisfy the same two conditions. The displayed
identity makes $g$ (like $\overline f$) $1$-Lipschitz, hence continuous. For
$\overline x = [(x_n)_n]$, the sequence $\iota_X(x_n)$ converges to $\overline x$ in
$\overline X$: indeed $\overline d(\iota_X(x_n),\overline x) = \lim_k d(x_n,x_k) \to 0$ as
$n\to\infty$, since $(x_n)_n$ is Cauchy. By continuity of $g$ and $f = g\circ\iota_X$,
\[g(\overline x) = \lim_n g(\iota_X(x_n)) = \lim_n f(x_n) = \overline f(\overline x),\]
using uniqueness of limits in $Y$ once more for the last equality. As $\overline x$ was
arbitrary, $g = \overline f$.
\end{exercisesolution}

\begin{exercisesolution}[Proof of \cref{ex:banach_hypotheses_sharp}]
% Generator: Claude Sonnet 5 (High)
\begin{enumerate}[label=\textbf{(\alph*)}]
  \item This is \cref{rem:banach_hypotheses}'s first bullet: $X := (0,1]$ is not complete,
    and $T(x) := x/2$ is a $\tfrac12$-contraction on $X$ with iterates marching towards the
    missing point $0$, hence no fixed point in $X$.
  \item Take $X := \mathbb{R}$ (complete, \cref{prop:Rn_coordinatewise_complete}) with the
    translation $T(x) := x+1$. Then $d(T(x),T(y)) = |(x+1)-(y+1)| = |x-y| = d(x,y)$, so $T$ is an
    isometry, and in particular Lipschitz with constant $L=1$ rather than $L<1$. Consequently,
    this does not contradict \cref{thm:banach_fixed_point}. A fixed point would need $x+1=x$,
    which is impossible, so $T$ has none.
\end{enumerate}
\end{exercisesolution}

\begin{exercisesolution}[Proof of \cref{ex:kobel_local_banach_fixed_point}]
% Generator: Claude Opus 5 (High)
Write $K := \overline{B}_r(x_0)$. The whole content of the exercise is that $f$ restricts to a
self-map of $K$, after which \cref{thm:banach_fixed_point} finishes the job.

\textbf{$K$ is complete.} A closed ball is a closed subset of $X$, and a closed subset of a
complete metric space is complete: a Cauchy sequence in $K$ is Cauchy in $X$, hence converges in
$X$, and its limit lies in $K$ because $K$ is closed. It is also non-empty, since $x_0 \in K$.

\textbf{$f(K) \subseteq K$.} Let $x \in K$, so $d(x,x_0) \le r$. Since $K \subseteq E$, both $x$
and $x_0$ lie in the domain of $f$, and the triangle inequality together with the contraction
estimate gives
\[d\big(f(x),x_0\big) \ \le\ d\big(f(x),f(x_0)\big) + d\big(f(x_0),x_0\big)
  \ \le\ L\,d(x,x_0) + (1-L)r \ \le\ Lr + (1-L)r \ =\ r,\]
so $f(x) \in K$. The hypothesis $d(f(x_0),x_0) \le (1-L)r$ is used exactly once, here, and it is
what makes the two terms add up to $r$ rather than to something larger.

\textbf{Conclusion.} The restriction $f|_K : K \to K$ is a contraction with the same constant $L$
on the non-empty complete metric space $K$, so \cref{thm:banach_fixed_point} supplies a unique
$\bar x \in K$ with $f(\bar x) = \bar x$. Since $K \subseteq E$, this $\bar x$ is a fixed point in
$E$, as required.

\textbf{Uniqueness comes for free, and in all of $E$.} \Cref{thm:banach_fixed_point} gives
uniqueness only within $K$, but the two-line argument used for the theorem itself needs nothing
beyond the contraction estimate, which holds on all of $E$: if $f(x)=x$ and $f(y)=y$ with
$x,y \in E$, then $d(x,y) = d(f(x),f(y)) \le L\,d(x,y)$, and $L<1$ forces $d(x,y)=0$. So $\bar x$
is the only fixed point of $f$ in $E$, even though the existence argument saw only the ball.
\end{exercisesolution}

\begin{exercisesolution}[Proof of \cref{ex:ai_contraction_cos}]
% Generator: Gemini 3.6 Flash (Medium)
\begin{enumerate}[label=\textbf{(\alph*)}]
  \item By the Mean Value Theorem, for any $x, y \in [0,1]$, there exists $\xi \in (0,1)$ such that $|f(x) - f(y)| = |f'(\xi)| |x-y| = \frac{1}{2}\sin(\xi) |x-y|$. Since $\sin(x)$ is strictly increasing on $[0,1]$, $\sup_{\xi \in [0,1]} \frac{1}{2}\sin(\xi) = \frac{1}{2}\sin(1) \approx 0.4207 < 1$. Thus $f$ is a contraction (\cref{def:contraction_mapping}) with $L = \frac{1}{2}\sin(1) < 1$.
  \item Since $[0,1]$ equipped with the Euclidean metric is a complete metric space (\cref{def:complete_metric_space}) and $f([0,1]) \subseteq [0,1]$, \cref{thm:banach_fixed_point} guarantees the existence and uniqueness of a fixed point $x^* \in [0,1]$ satisfying $f(x^*) = x^*$.
\end{enumerate}
\end{exercisesolution}

\begin{exercisesolution}[Solution to \cref{ex:3.5}]
% Generator: Claude Opus 5 (High)
\textbf{(a), (c) and (e) are true; (b) and (d) are false.}

\begin{enumerate}[label=\textbf{(\alph*)}]
  \item \textbf{True,} and this is the classic illustration of \cref{thm:banach_fixed_point}. The
    map of Zurich lying on the desk defines a function from the (compact, hence complete) region
    of the city depicted onto a smaller region of the desk, and that function is a similarity: a
    rotation and translation composed with a scaling by the map's ratio $L < 1$. It is therefore
    a contraction of a complete metric space into itself, and so has exactly one fixed point: one
    spot on the map lying directly above the point of the city it represents.
  \item \textbf{False.} The hypothesis says nothing about $f$ mapping its domain into itself, and
    the codomain $[0,2]$ is strictly larger than the domain $[0,1]$. Take
    \[f(x) := \tfrac{x}{2} + \tfrac32, \qquad f : [0,1] \to [\tfrac32,2] \subseteq [0,2],\]
    which is $\tfrac12$-Lipschitz. A fixed point would need $x = \tfrac x2 + \tfrac32$, i.e.\
    $x=3$, which is not in $[0,1]$.
  \item \textbf{True.} Now the inclusion runs the right way: $f$ maps $[0,2]$ into
    $[0,1] \subseteq [0,2]$, so $f$ is a self-map of the complete space $[0,2]$
    (closed and bounded in $\mathbb{R}$, hence compact, hence complete by
    \cref{ex:compact_finite_complete}), and it is a contraction with $L = \tfrac12 < 1$.
    \Cref{thm:banach_fixed_point} applies directly.
  \item \textbf{False,} and this is the instructive one: the domain is \emph{disconnected}, and
    the condition $|f'| < 1$ is purely local, so it cannot see across the gap. Define
    \[f(x) := \begin{cases} \tfrac52 & x \in [0,1], \\ \tfrac12 & x \in [2,3].\end{cases}\]
    Then $f$ is continuously differentiable with $f' \equiv 0$, it maps
    $[0,1]\cup[2,3]$ into itself, and it has no fixed point, because each of the two pieces is
    sent into the \emph{other} one. Compare the local-to-global reading of
    \cref{rem:local_to_global}: a pointwise derivative bound upgrades to a global contraction
    only along paths that stay inside the domain, and here there are none.
  \item \textbf{True,} but \emph{not} by the fixed point theorem, and that is the point of the
    star. The hypothesis $|f'(x)| < 1$ gives no uniform $L < 1$, since the supremum of $|f'|$ may
    well equal $1$, so \cref{thm:banach_fixed_point} is unavailable. Existence follows from the
    intermediate value theorem instead. Put $g(x) := f(x) - x$, which is continuous on $[0,1]$.
    Since $f$ takes values in $[0,1]$,
    \[g(0) = f(0) \geq 0 \qquad\text{and}\qquad g(1) = f(1) - 1 \leq 0,\]
    so $g$ vanishes somewhere in $[0,1]$ by \cref{cor:intermediate_value_theorem}, and that zero
    is a fixed point of $f$. Note that the derivative hypothesis was never used: \emph{every}
    continuous self-map of $[0,1]$ has a fixed point. What $|f'| < 1$ buys is
    \textbf{uniqueness}, by the same two-line argument as in the proof of uniqueness for
    \cref{thm:banach_fixed_point} combined with the mean value theorem.
\end{enumerate}
\end{exercisesolution}

\begin{exercisesolution}[Solution to \cref{ex:bfpt_applicability_table}]
% Generator: Claude Opus 5 (High)
Recall that \cref{thm:banach_fixed_point} needs three things at once: $X$ complete, $T$ mapping
$X$ into $X$, and a Lipschitz constant $\lambda < 1$ that is \emph{actually attained as a bound},
not merely approached.
\begin{itemize}
  \item \textbf{$T(x) = \sin(x)$ on $\mathbb{R}$. Does not apply; fixed points $\{0\}$.}
    $\mathbb{R}$ is complete and $T$ maps it into itself, but
    $\lambda = \sup_{x}\lvert\cos x\rvert = 1$, attained at $x=0$. So the theorem is unavailable.
    A fixed point exists all the same: $\sin x = x$ holds at $x=0$ and nowhere else, since
    $\lvert\sin x\rvert < \lvert x\rvert$ for $x \neq 0$. This is the first lesson of the table:
    the theorem is \emph{sufficient}, never necessary.
  \item \textbf{$T(x) = \tfrac{9}{10}\cos(x)$ on $\mathbb{R}$. Applies; one fixed point.}
    Here $\lvert T'(x)\rvert = \tfrac{9}{10}\lvert\sin x\rvert \leq \tfrac{9}{10} < 1$, so by the
    mean value theorem $T$ is a $\tfrac{9}{10}$-contraction; $\mathbb{R}$ is complete and
    $T(\mathbb{R}) \subseteq [-\tfrac9{10},\tfrac9{10}] \subseteq \mathbb{R}$. All three
    hypotheses hold, so there is exactly one fixed point. It has no closed form; iterating from
    any starting value gives $x^* \approx 0.693$.
  \item \textbf{$T(x) = x^2$ on $\mathbb{R}$. Does not apply; fixed points $\{0,1\}$.}
    $T$ is not Lipschitz at all on $\mathbb{R}$, since $T'(x) = 2x$ is unbounded. That it cannot
    apply is visible without any computation: $x^2 = x$ has the two solutions $0$ and $1$, and
    the theorem asserts \emph{uniqueness}, so any set of hypotheses implying it must fail here.
  \item \textbf{$T(x) = \tfrac12x^2$ on $(-1,1)$. Does not apply; fixed points $\{0\}$.}
    Two hypotheses fail together. The domain $(-1,1)$ is not complete, and
    $\lambda = \sup_{\lvert x\rvert<1}\lvert T'(x)\rvert = \sup\lvert x\rvert = 1$. This supremum
    is not attained, but a supremum of $1$ is still not $<1$, which is exactly the trap of
    \cref{rem:banach_hypotheses}. The map does send $(-1,1)$ into $[0,\tfrac12) \subseteq
    (-1,1)$. Solving $\tfrac12x^2 = x$ gives $x \in \{0,2\}$, of which only $0$ lies in the
    domain, so a unique fixed point exists regardless.
  \item \textbf{$T(x) = \tfrac12\big(x+\tfrac2x\big)$ on $\mathbb{Q}\cap(1,2)$. Does not apply;
    \emph{no} fixed points.} This row is the interesting one, and it is the Babylonian iteration
    for $\sqrt2$. Everything works except completeness:
    \begin{itemize}
      \item $T'(x) = \tfrac12 - \tfrac1{x^2}$, so on $(1,2)$ we have
        $-\tfrac12 \leq T'(x) \leq \tfrac14$ and hence $\lambda = \tfrac12 < 1$: a genuine
        contraction.
      \item $T$ maps $X$ into $X$: it sends rationals to rationals, and on $(1,2)$ its minimum
        is $T(\sqrt2) = \sqrt2 \approx 1.414$ while $T(1) = T(2) = \tfrac32$, so
        $T\big((1,2)\big) \subseteq [\sqrt2,\tfrac32] \subseteq (1,2)$.
      \item But $\mathbb{Q}\cap(1,2)$ is \textbf{not complete}.
    \end{itemize}
    And there is indeed no fixed point: $x = \tfrac12(x+\tfrac2x)$ forces $x^2 = 2$, so
    $x = \sqrt2 \notin \mathbb{Q}$. The iterates are the familiar rational approximations
    $\tfrac32, \tfrac{17}{12}, \tfrac{577}{408}, \dots$, which form a Cauchy sequence in $X$
    converging to a point that $X$ does not contain.

    This is worth contrasting with \cref{rem:banach_hypotheses}, where completeness failed on
    $(0,1]$ with the missing point sitting at an endpoint one is already suspicious of. Here the
    domain is an interval of rationals, the hole is invisible, and the contraction is a
    thoroughly practical algorithm. Completeness is not a technicality.
\end{itemize}
\end{exercisesolution}

\begin{exercisesolution}[Solution of \cref{ex:incomplete_metric_space_example}]
% Generator: Claude Opus 5 (High)
Take $X := (0,1]$ with the Euclidean distance $d(x,y) := \lvert x-y\rvert$ inherited from
$\mathbb{R}$. The sequence $x_k := 1/k$ lies in $X$ and is Cauchy, since $\lvert x_k - x_l\rvert \le
\max\{1/k, 1/l\} \to 0$. But it does not converge in $X$: its only candidate limit in $\mathbb{R}$
is $0$, and $0 \notin X$, so no point of $X$ can be its limit. Hence $(X,d)$ is not complete.

The same mechanism gives the other standard answers. Taking $X := \mathbb{Q}$ with
$d(x,y) := \lvert x-y\rvert$, any rational sequence converging to $\sqrt2$ is Cauchy without a limit
in $X$; taking $X := \mathbb{R} \setminus \{0\}$ works for the same reason as $(0,1]$.
\begin{ainote}
What is being exploited each time is that completeness is \emph{not} a property of the distance
alone: $(0,1]$ and $[0,1]$ carry the same formula for $d$, and one is complete while the other is
not. A Cauchy sequence is a sequence that is trying to converge, and the question is only whether
the space has kept the point it is aiming at. This is exactly the gap that the completion of a
metric space (\cref{sec:completion_of_a_metric_space}) is built to fill.
\end{ainote}
\end{exercisesolution}

\begin{exercisesolution}[Solution of \cref{ex:kobel_cauchy_single_choice}]
% Generator: Claude Opus 5 (High)
Only \textbf{(d)} is correct.
\begin{enumerate}[label=\textbf{(\alph*)}]
  \item \textbf{False.} A Cauchy sequence converges precisely when the ambient space has kept the
    point it is aiming at, and that is the definition of completeness
    (\cref{def:complete_metric_space}), not a theorem about Cauchy sequences.
    \Cref{ex:incomplete_metric_space_example} is a counterexample in one line.
  \item \textbf{False}, and this is the option the question is built around. Closedness is not a
    property a space has, only a property a subset has \emph{relative to} something larger: every
    metric space is closed in itself, so reading (b) as a hypothesis on $X$ alone makes it say
    nothing at all, and $X := (0,1]$ is again a counterexample. What is true is the relative
    statement: a closed subset of a \emph{complete} space is complete, and the ambient
    completeness is the part (b) omits.
  \item \textbf{False.} An interval need not be complete either. On $X := (0,1)$ the sequence
    $a_i := 1/i$ is Cauchy with no limit in $X$. Only the closed bounded intervals, and the closed
    unbounded ones, are complete.
  \item \textbf{True.} A compact metric space is sequentially compact
    (\cref{thm:sequential_compactness}), so $(a_i)$ has a convergent subsequence, and a Cauchy
    sequence with a convergent subsequence converges
    (\cref{ex:cauchy_elementary_facts}\textbf{(c)}). Hence every compact metric space is complete,
    which is \cref{ex:compact_finite_complete}.
\end{enumerate}
\end{exercisesolution}

\begin{exercisesolution}[Solution of \cref{ex:banach_hypotheses_true_false}]
% Generator: Claude Opus 5 (High)
\begin{enumerate}[label=\textbf{(\alph*)}]
  \item \textbf{True.} For $x,y \in [0,1]$,
    $\lvert x^2 - y^2\rvert = \lvert x+y\rvert\,\lvert x-y\rvert \le 2\lvert x-y\rvert$,
    since $x + y \le 2$ on $[0,1]$. So $2$ is a Lipschitz constant. It is even the optimal one, as
    $\lvert f'(x)\rvert = 2x$ approaches $2$ at $x = 1$.
  \item \textbf{False.} The space $X = [0,1]$ is complete, and $f$ does map $X$ into $X$; what fails
    is the contraction hypothesis. A contraction needs a Lipschitz constant \emph{strictly less
    than $1$}, and by \textbf{(a)} the best available constant here is $2$. Concretely, taking
    $x = 1$ and $y = 1 - \varepsilon$ gives
    $\lvert f(x) - f(y)\rvert = (2-\varepsilon)\varepsilon > \lvert x - y\rvert$ for small
    $\varepsilon > 0$: near $x = 1$ the map \emph{expands} distances.
  \item \textbf{True} --- and note the plural. Solving $x^2 = x$ on $[0,1]$ gives $x = 0$ and
    $x = 1$, so $f$ has exactly two fixed points.
\end{enumerate}
\begin{remark}
Parts \textbf{(b)} and \textbf{(c)} together are the point of the exercise. The Banach fixed point
theorem (\cref{thm:banach_fixed_point}) is a sufficient condition, not a necessary one: its
hypotheses fail here, and fixed points exist anyway. What the failure does cost is the
\emph{uniqueness} half of the conclusion, and sure enough there are two of them. Reading \qt{the
hypotheses fail} as \qt{there is no fixed point} is the standard mistake.
\end{remark}
\end{exercisesolution}

\begin{exercisesolution}[Proof of \cref{ex:banach_fixed_point_linear}]
% Generator: Gemini 3.6 Flash (High)
% Correction: Claude Opus 5 (High) --- rewritten for R^n. Flash's version worked in an abstract
% metric space and wrote d(f(x)/lambda, f(y)/lambda) = d(f(x),f(y))/lambda, which is meaningless
% there and false for a general metric even when scaling does make sense. It is the homogeneity of
% the Euclidean norm that is doing the work, so the norm has to be visible in the computation.
% Confirmed against WS22_Lösung.pdf p. 16, which runs the identical norm computation and likewise
% names the threshold lambda > L. The official solution agrees line for line.
\begin{enumerate}[label=\textbf{(\alph*)}]
  \item Fix $x,y \in \mathbb{R}^n$. Since the Euclidean norm is homogeneous and $f$ is
    $L$-Lipschitz,
\[
\lvert g(x) - g(y)\rvert = \left\lvert \frac{f(x) - f(y)}{\lambda} \right\rvert
  = \frac{1}{\lambda}\,\lvert f(x) - f(y)\rvert \leq \frac{L}{\lambda}\,\lvert x - y\rvert .
\]
Set $k := L/\lambda$. Since $\lambda > L > 0$ we have $k < 1$, and therefore $g$ is a contraction
with contraction constant $k$.

  \item The space $\mathbb{R}^n$ is complete, so the Banach Fixed Point Theorem
    (\cref{thm:banach_fixed_point}) applies to $g$ and yields a unique $x \in \mathbb{R}^n$ with
    $g(x) = x$. Multiplying $\lambda^{-1} f(x) = x$ by $\lambda$ turns this into $f(x) = \lambda x$,
    and the two equations have exactly the same solutions. Hence $f(x) = \lambda x$ has a unique
    solution.
\end{enumerate}
\begin{ainote}
The hypothesis is what makes this work, and it is worth naming: a Lipschitz $f$ need not be a
contraction, and need have no fixed point at all (translation $x \mapsto x + v$ is $1$-Lipschitz).
Dividing by $\lambda > L$ buys the missing factor. This is also why the original exam question says
\qt{for all sufficiently large $\lambda$} rather than naming a threshold: any $\lambda > L$ works,
but $L$ is not part of the data the question hands you.
\end{ainote}
\end{exercisesolution}

\begin{exercisesolution}[Proof of \cref{ex:examHS24_TF_banach}]
% Generator: Gemini 3.1 Pro (High)
\begin{enumerate}[label=\textbf{(\alph*)}]
  \item \textbf{True:} The derivative is $f'(x) = -\sin(x)$. Since $\lvert f'(x)\rvert = \lvert\sin x\rvert \le 1$ for all $x \in \mathbb{R}$, the mean value theorem gives $\lvert f(x) - f(y)\rvert \le 1 \cdot \lvert x-y\rvert$, so $f$ has Lipschitz constant $1$.
  \item \textbf{False:} The Banach fixed point theorem requires a \emph{strict} contraction, that is, a Lipschitz constant $\lambda < 1$. Here, the tightest possible constant is $\sup_{x \in [-\pi,\pi]} \lvert f'(x)\rvert = 1$ (attained at $x = \pm \pi/2$), so the map is non-expansive but not strictly contractive.
  \item \textbf{True:} The function $f(x) = \cos x$ is continuous and maps $[-\pi,\pi]$ into $[-1,1] \subseteq [-\pi,\pi]$. Since $g(x) := \cos x - x$ has $g(-\pi) = \pi - 1 > 0$ and $g(\pi) = -1 - \pi < 0$, the intermediate value theorem supplies a root, which is a fixed point of $f$.
\end{enumerate}
\end{exercisesolution}


